\documentclass[11pt]{article}
% ===================== PREAMBUŁA =====================
\usepackage[polish]{babel}
\usepackage[T1]{fontenc}
\usepackage[utf8]{inputenc}
\usepackage{geometry}
\usepackage{graphicx}
\usepackage{booktabs}
\usepackage{amsmath}
\usepackage{tikz}
\usepackage{caption}
\usepackage{subcaption}
\usepackage{multirow}
\geometry{margin=2.5cm}

% --- biblioteki TikZ do rysunkow logiczno-czasowych (rysunki inline, bez plikow PNG) ---
\usetikzlibrary{shapes.gates.logic.US}
\usetikzlibrary{arrows.meta}
\usetikzlibrary{positioning}
\usetikzlibrary{calc}
% ===================== RYSUNKI (inline TikZ) =====================
% =====================================================================
%  RYSUNKI LOGICZNE (inline TikZ) — wspolny plik dolaczany przez \input
%  Wymaga w preamble:
%  \usetikzlibrary{shapes.gates.logic.US, arrows.meta, positioning, calc}
% =====================================================================

% ---------------------------------------------------------------------
% RYSUNEK A — SCHEMAT BLOKOWY SUMATORA 2-BITOWEGO
% ---------------------------------------------------------------------
\newcommand{\RysSchematBlokowy}{%
\begin{tikzpicture}[
  font=\small,
  line width=0.7pt,
  >=Latex,
  blok/.style={draw, thick, minimum width=32mm, minimum height=13mm, align=center},
  dotTikZ/.style={draw, circle, inner sep=1.1pt, fill=black}
]
  \node[blok, fill=gray!10] (hs) at (0,0)   {półsumator\\[-2pt]$S_0=A_0\oplus B_0$\\[-2pt]$P_0=A_0\cdot B_0$};
  \node[blok, fill=gray!10] (fs) at (7.0,0) {pełny sumator\\[-2pt]$S_1=A_1\oplus B_1\oplus P_0$\\[-2pt]$P_1$};
  \node[blok, minimum width=24mm, minimum height=9mm, fill=gray!10] (wys) at (7.0,2.9) {$S_2=P_1$};

  \draw[->] (-4.6,0.8)  node[left] {$A_0$} -- (hs.160);
  \draw[->] (-4.6,-0.8) node[left] {$B_0$} -- (hs.200);
  \draw[->] (-4.6,2.1)  node[left] {$A_1$} -- (3.6,2.1) -- (fs.150);
  \draw[->] (-4.6,-2.1) node[left] {$B_1$} -- (5.4,-2.1) -- (fs.210);
  \draw[->] (hs.east) -- ++(0.5,0) coordinate(pd) node[dotTikZ] {} node[above=1pt] {$P_0$};
  \draw[->] (pd) -- (fs.205);
  \draw[->] (hs.north) |- (9.6,-1.6) node[dotTikZ] {} |- ($(fs.east)+(0.35,-1.0)$);
  \draw[->] ($(hs.east)+(0.25,0.45)$) node[dotTikZ] {} -- ++(0.4,0) |- (9.6,-2.6) node[right] {$S_0$};
  \draw[->] (fs.north) -- (wys.south);
  \draw[->] (wys.east) -- ++(2.0,0) node[right] {$S_2$};
  \draw[->] (fs.east) -- ++(2.0,0) node[right] {$S_1$};
\end{tikzpicture}%
}

% ---------------------------------------------------------------------
% RYSUNEK B — UKŁAD HIERARCHICZNY: DWA POZIOMY SUMACYJNE
% ---------------------------------------------------------------------
\newcommand{\RysSchematBlokowyDwa}{%
\begin{tikzpicture}[
  font=\small,
  line width=0.7pt,
  >=Latex,
  blok/.style={draw, thick, minimum width=34mm, minimum height=12mm, align=center}
]
  \node[blok, fill=gray!10] (p1) at (0,1.6)   {poziom młodszy\\[-2pt]$S_0,\ P_0$};
  \node[blok, fill=gray!10] (p2) at (6.8,1.6) {poziom starszy\\[-2pt]$S_1,\ P_1$};
  \node[blok, minimum width=20mm, minimum height=8mm, fill=gray!10] (p3) at (6.8,4.1) {$S_2$};

  \draw[->] (-4.8,2.2) node[left] {$A_0,\ B_0$} -- (p1.west);
  \draw[->] (-4.8,1.0) node[left] {$A_1,\ B_1$} -- (-2.8,1.0) -- (-2.8,4.1) -| (p2.west);
  \draw[->] (p1.east) -- node[above] {$P_0$} (p2.west);
  \draw[->] (p1.north) |- (9.4,-2.2) node[right] {$S_0$};
  \draw[->] (p2.north) -- (p3.south);
  \draw[->] (p2.east) -- ++(2.2,0) node[right] {$S_1$};
  \draw[->] (p3.east) -- ++(2.2,0) node[right] {$S_2$};
\end{tikzpicture}%
}

% ---------------------------------------------------------------------
% RYSUNEK C — PEŁNY SCHEMAT LOGICZNY NA BRAMKACH NOR (17 BRAMEK)
%  S0 = XOR(A0,B0);  P0 = A0*B0;
%  X  = XOR(A1,B1);  S1 = XOR(X,P0);
%  S2 = P1 = NOR( NOR(A1,B1), NOR(X,P0) )
%  XOR(u,v) z 4 NOR: N1=NOR(u,v), N2=NOR(u,~N1), N3=NOR(v,~N1), out=NOR(N2,N3)
% ---------------------------------------------------------------------
\newcommand{\RysSchematNOR}{%
\begin{tikzpicture}[
  font=\small,
  line width=0.6pt,
  >=Latex,
  norI2/.style={nor gate US, draw, logic gate inputs=nn, minimum width=11mm},
  dotTikZ/.style={draw, circle, inner sep=1.1pt, fill=black},
  lbl/.style={font=\footnotesize}
]
  % ================= POZIOM MLodszy: S0, P0 =================
  % --- bramki ---
  \node[norI2] (iA0)   at (0,1.2)    {};
  \node[norI2] (iB0)   at (0,-0.4)   {};
  \node[norI2] (gXOR1) at (2.8,0.4)  {};
  \node[norI2] (gXOR2) at (2.8,-1.6) {};
  \node[norI2] (gS0)   at (5.6,-0.6) {};
  \node[norI2] (gP0)   at (5.6,-2.6) {};
  % --- wejscia A0, B0 + inwertery ---
  \draw (iA0.input 1) -- ++(-0.6,0) coordinate(A0) node[left] {$A_0$};
  \draw (iB0.input 1) -- ++(-0.6,0) coordinate(B0) node[left] {$B_0$};
  \draw ($(A0)+(-0.6,0)$) node[dotTikZ] {} coordinate(A0b);
  \draw ($(B0)+(-0.6,0)$) node[dotTikZ] {} coordinate(B0b);
  \draw (iA0.input 1) -- (iA0.input 2);
  \draw (iB0.input 1) -- (iB0.input 2);
  % --- XOR(A0,B0) ---
  \draw (A0b) |- ($(gXOR1.input 1)+(-0.3,0)$) -- (gXOR1.input 1);
  \draw (B0)  |- ($(gXOR1.input 2)+(-0.3,0)$) -- (gXOR1.input 2);
  \draw (iA0.output) -- ++(0.3,0) coordinate(v1);
  \draw (iB0.output) -- ++(0.3,0) coordinate(v2);
  \draw (v1) |- ($(gXOR2.input 1)+(-0.3,0)$) -- (gXOR2.input 1);
  \draw (v2) |- ($(gXOR2.input 2)+(-0.3,0)$) -- (gXOR2.input 2);
  % --- NOR = OR dla S0 oraz P0 ---
  \draw (gXOR1.output) -- ++(0.3,0) coordinate(c1) node[dotTikZ] {};
  \draw (gXOR2.output) -- ++(0.3,0) coordinate(c2);
  \draw (c1) |- ($(gS0.input 1)+(-0.3,0)$) -- (gS0.input 1);
  \draw (c1) |- ($(gP0.input 1)+(-0.3,0)$) -- (gP0.input 1);
  \draw (c2) |- ($(gS0.input 2)+(-0.3,0)$) -- (gS0.input 2);
  \draw (gS0.output) -- ++(0.5,0) node[dotTikZ] {} coordinate(S0b) node[right, lbl] {$S_0$};
  \draw (S0b) |- (gP0.input 2);
  \draw (gP0.output) -- ++(0.6,0) node[dotTikZ] {} coordinate(P0) node[below=3pt, lbl] {$P_0$};

  % ================= POZIOM STARSZY: S1, S2=P1 =================
  % --- bramki ---
  \node[norI2] (iA1)   at (0,8.2)    {};
  \node[norI2] (iB1)   at (0,6.6)    {};
  \node[norI2] (hXOR1) at (2.8,7.4)  {};
  \node[norI2] (hXOR2) at (2.8,5.4)  {};
  \node[norI2] (hX)    at (5.6,6.4)  {};
  \node[norI2] (iX)    at (5.6,4.2)  {};
  \node[norI2] (hXP)   at (8.6,5.3)  {};
  \node[norI2] (iP0h)  at (8.6,2.9)  {};
  \node[norI2] (hXP2)  at (11.4,4.1) {};
  \node[norI2] (hS1)   at (14.2,4.7) {};
  \node[norI2] (hP1)   at (14.2,7.0) {};
  % --- wejscia A1, B1 + inwertery ---
  \draw (iA1.input 1) -- ++(-0.6,0) coordinate(A1) node[left] {$A_1$};
  \draw (iB1.input 1) -- ++(-0.6,0) coordinate(B1) node[left] {$B_1$};
  \draw ($(A1)+(-0.6,0)$) node[dotTikZ] {} coordinate(A1b);
  \draw ($(B1)+(-0.6,0)$) node[dotTikZ] {} coordinate(B1b);
  \draw (iA1.input 1) -- (iA1.input 2);
  \draw (iB1.input 1) -- (iB1.input 2);
  % --- X = XOR(A1,B1) ---
  \draw (A1b) |- ($(hXOR1.input 1)+(-0.3,0)$) -- (hXOR1.input 1);
  \draw (B1b) |- ($(hXOR1.input 2)+(-0.3,0)$) -- (hXOR1.input 2);
  \draw (iA1.output) -- ++(0.3,0) coordinate(w1);
  \draw (iB1.output) -- ++(0.3,0) coordinate(w2);
  \draw (w1) |- ($(hXOR2.input 1)+(-0.3,0)$) -- (hXOR2.input 1);
  \draw (w2) |- ($(hXOR2.input 2)+(-0.3,0)$) -- (hXOR2.input 2);
  \draw (hXOR1.output) -- ++(0.3,0) coordinate(d1) node[dotTikZ] {};
  \draw (hXOR2.output) -- ++(0.3,0) coordinate(d2);
  \draw (d1) |- ($(hX.input 1)+(-0.3,0)$) -- (hX.input 1);
  \draw (d2) |- ($(hX.input 2)+(-0.3,0)$) -- (hX.input 2);
  % --- inwerter X -> notX ---
  \draw (hX.output) -- ++(0.25,0) node[dotTikZ] {} coordinate(Xb);
  \draw (Xb) |- (iX.input 1);
  \draw (iX.input 1) -- (iX.input 2);
  \draw (iX.output) -- ++(0.3,0) coordinate(e1);
  % --- S1 = NOR( NOR(X,P0), NOR(notX, notP0) ) ---
  \draw (Xb) |- ($(hXP.input 1)+(-0.3,0)$) -- (hXP.input 1);
  \draw (e1) |- ($(hXP2.input 1)+(-0.3,0)$) -- (hXP2.input 1);
  \draw (P0) |- (iP0h.input 1);
  \draw (iP0h.input 1) -- (iP0h.input 2);
  \draw (iP0h.output) -- ++(0.3,0) coordinate(e2);
  \draw (e2) |- ($(hXP2.input 2)+(-0.3,0)$) -- (hXP2.input 2);
  \draw (P0) |- ($(hXP.input 2)+(-0.3,0)$) -- (hXP.input 2);
  \draw (hXP.output) -- ++(0.3,0) coordinate(f1) node[dotTikZ] {};
  \draw (hXP2.output) -- ++(0.3,0) coordinate(f2);
  \draw (f1) |- ($(hS1.input 1)+(-0.3,0)$) -- (hS1.input 1);
  \draw (f2) |- ($(hS1.input 2)+(-0.3,0)$) -- (hS1.input 2);
  \draw (hS1.output) -- ++(0.5,0) node[right, lbl] {$S_1$};
  % --- S2 = P1 = NOR( NOR(A1,B1), NOR(X,P0) ) ---
  \draw (d1) |- ($(hP1.input 1)+(-0.3,0)$) -- (hP1.input 1);
  \draw (f1) |- ($(hP1.input 2)+(-0.3,0)$) -- (hP1.input 2);
  \draw (hP1.output) -- ++(0.7,0) node[right, lbl] {$S_2=P_1$};
  % --- etykiety poziomow ---
  \node[lbl, gray] at (2.2,-3.9) {poziom młodszy: $S_0$, $P_0$};
  \node[lbl, gray] at (7.2,9.4)  {poziom starszy: $S_1$, $S_2=P_1$};
\end{tikzpicture}%
}

% alias dla punktu 3 (schemat zminimalizowany = ta sama siec NOR)
\newcommand{\RysSchematMinimalny}{\RysSchematNOR}

% ---------------------------------------------------------------------
% RYSUNEK D — PRZEBIEGI CZASOWE (16 KOMBINACJI WEJŚĆ, jak w tablicy prawdy)
% ---------------------------------------------------------------------
\newcommand{\RysPrzebiegi}{%
\begin{tikzpicture}[
  font=\footnotesize,
  >=Latex,
  sig/.style={line width=1.0pt},
  hi/.style={line width=1.0pt},
  gr/.style={gray!45, very thin}
]
  \def\dx{0.62}
  % siatka pionowa t = 0..16
  \foreach \x in {0,1,...,16} {
    \draw[gr] (2+\x*\dx, 0.1) -- (2+\x*\dx, 8.4);
  }
  % ---------- A1: 0 dla t<8, 1 dla t>=8 ----------
  \draw[sig] (2,7.9) node[left] {$A_1$} -- (2+16*\dx,7.9);
  \draw[hi]  (2,7.2) -- (2+8*\dx,7.2);
  \draw[hi]  (2+8*\dx,7.2) -- (2+8*\dx,7.9) -- (2+16*\dx,7.9);
  % ---------- A0: 0,0,1,1 cyklicznie ----------
  \draw[sig] (2,6.1) node[left] {$A_0$} -- (2+16*\dx,6.1);
  \foreach \a/\b in {2/4,6/8,10/12,14/16} {
    \draw[hi] (2+\a*\dx,5.4) -- (2+\b*\dx,5.4);
    \draw[hi] (2+\a*\dx,5.4) -- (2+\a*\dx,6.1);
    \draw[hi] (2+\b*\dx,5.4) -- (2+\b*\dx,6.1);
  }
  % ---------- B1: 0,0,1,1 cyklicznie (wolniej) ----------
  \draw[sig] (2,4.3) node[left] {$B_1$} -- (2+16*\dx,4.3);
  \foreach \a/\b in {4/8,12/16} {
    \draw[hi] (2+\a*\dx,3.6) -- (2+\b*\dx,3.6);
    \draw[hi] (2+\a*\dx,3.6) -- (2+\a*\dx,4.3);
    \draw[hi] (2+\b*\dx,3.6) -- (2+\b*\dx,4.3);
  }
  % ---------- B0: 0,1,0,1 cyklicznie (najszybciej) ----------
  \draw[sig] (2,2.5) node[left] {$B_0$} -- (2+16*\dx,2.5);
  \foreach \a in {1,3,...,15} {
    \pgfmathsetmacro\bp{\a+1}
    \draw[hi] (2+\a*\dx,1.8) -- (2+\bp*\dx,1.8);
    \draw[hi] (2+\a*\dx,1.8) -- (2+\a*\dx,2.5);
    \draw[hi] (2+\bp*\dx,1.8) -- (2+\bp*\dx,2.5);
  }
  % ---------- S2: 1 gdy suma >= 4 (t = 4..7, 10..11, 13..16) ----------
  \draw[sig] (2,0.7) node[left] {$S_2$} -- (2+16*\dx,0.7);
  \foreach \a/\b in {4/8,10/12,13/16} {
    \draw[hi] (2+\a*\dx,0.0) -- (2+\b*\dx,0.0);
    \draw[hi] (2+\a*\dx,0.0) -- (2+\a*\dx,0.7);
    \draw[hi] (2+\b*\dx,0.0) -- (2+\b*\dx,0.7);
  }
  % ---------- podpisy stanow A1A0 / B1B0 pod osia ----------
  \foreach \i [count=\j from 0] in {00,00,01,01,10,10,11,11,00,00,01,01,10,10,11,11} {
    \node[rotate=90, anchor=east, font=\tiny] at (2+\j*\dx+0.5*\dx, -0.25) {\i};
  }
  \foreach \i [count=\j from 0] in {00,01,10,11,00,01,10,11,00,01,10,11,00,01,10,11} {
    \node[rotate=90, anchor=east, font=\tiny] at (2+\j*\dx+0.5*\dx, -1.15) {\i};
  }
  \node[font=\tiny, anchor=north] at (2+8*\dx, -2.35) {stany $A_1A_0$ (górny rząd) i $B_1B_0$ (dolny rząd) zgodne z tablicą prawdy};
\end{tikzpicture}%
}
\title{Projekt układu sumy liczby 2-bitowej z 3-bitową na bramkach NOR}
\author{}
\date{}

\begin{document}

\maketitle
\tableofcontents
\newpage

\section{Wstęp teoretyczny}

\subsection{Definicja sumatora cyfrowego}

Sumator cyfrowy jest układem kombinacyjnym służącym do wykonywania dodawania liczb zapisanych w postaci binarnej. Jego zadaniem jest pozyskanie wyniku sumy oraz ewentualnego przeniesienia na podstawie wejściowych słów bitowych. W zależności od zaprojektowanej architektury sumator może realizować dodawanie dwóch liczb bez przeniesienia, z jednym przeniesieniem właściwym lub z mechanizmem przeniesienia wyprowadzonego na wyjście.

W niniejszym projekcie analizowany jest układ dodający dwie liczby dwubitowe, dla których wynik sumy może przyjąć wartość 3-bitową. Takie zadanie wymaga uwzględnienia możliwości wystąpienia przeniesienia z najstarszego bitu wyniku, co w naturalny sposób prowadzi do konieczności wyznaczenia trzech bitów wyniku: dwóch bitów sumy oraz jednego bitu przeniesienia.

Dla przyjętego rozwiązania:
\begin{itemize}
    \item liczba pierwsza przedstawiana jest bitami \(A_1, A_0\) (gdzie \(A_1\) jest bitem starszym, a \(A_0\) bity młodszym),
    \item liczba druga przedstawiana jest bitami \(B_1, B_0\),
    \item wynik sumy 3-bitowego oznaczany jest jako \(S_2 S_1 S_0\),
    \item bity \(S_1, S_0\) odpowiadają za wynik dodawania, a bit \(S_2\) odpowiada za najstarsze przeniesienie.
\end{itemize}

Wyjście układu oznaczone jako \(S_1, S_0, S_2\) przyjmuje wartości zależne od sumy wejściowych liczb dwu-bitowych. Najstarszy bit wyniku \(S_2\) przyjmuje wartość logiczną 1 wtedy, gdy suma liczb wejściowych jest większa lub równa 4; w pozostałych przypadkach przyjmuje wartość 0.

Suma dwóch liczb dwubitowych może być przedstawiona jako:
\begin{equation}
A + B = S
\end{equation}
gdzie:
\begin{itemize}
    \item \(A = 2A_1 + A_0\),
    \item \(B = 2B_1 + B_0\),
    \item \(S = 4S_2 + 2S_1 + S_0\).
\end{itemize}

\subsection{Zasada działania sumatora 2-bitowego z 3-bitową sumą}

Dodawanie dwu-bitowych liczb binarnych polega na sumowaniu odpowiadających sobie bitów z uwzględnieniem przeniesienia z dolnego poziomu na wyższy. W najprostszej formie można to realizować przy użyciu podstawowych bramek logicznych, w tym bramek XOR, AND oraz OR. Jednak w niniejszym projekcie przyjmuje się ograniczenie do realizacji wyłącznie na bramkach NOR, co wprowadza dodatkowe wymagania w zakresie minimalizacji i zamiany funkcji logicznych na postać przyjazną dla bramek NOR.

Dla liczb dwubitowych:
\begin{center}
$A = (A_1 A_0),\quad B = (B_1 B_0)$
\end{center}
należy wykonać dodawanie pozycyjne, począwszy od bitu najmłodszej pozycji, a następnie przenieść ewentualne przeniesienie na bit starszy. W rezultacie otrzymujemy:
\begin{itemize}
    \item \(S_0\) -- wynik sumy bitów najmłodszych \(A_0\) i \(B_0\) z uwzględnieniem przeniesienia 0 na tym poziomie,
    \item \(S_1\) -- wynik sumy bitów starszych \(A_1\) i \(B_1\) oraz przeniesienia z niższego poziomu,
    \item \(S_2\) -- przeniesienie z poziomu bitów starszych.
\end{itemize}

W praktyce realizacji na bramkach NOR każda z tych funkcji może być przekształcana do postaci sumy iloczynów lub iloczynu sum, a następnie minimalizowana w taki sposób, aby struktura miała możliwość realizacji przy użyciu wyłącznie bramek NOR.

Można to opisać za pomocą schematu działania układu w następujący sposób:
\begin{enumerate}
    \item poziom młodszy: sumowanie \(A_0, B_0\) z uzyskaniem \(S_0\) oraz przeniesienia \(P_0\),
    \item poziom starszy: sumowanie \(A_1, B_1\) oraz przeniesienia \(P_0\) z uzyskaniem \(S_1\) i przeniesienia \(P_1\),
    \item wyjście najstarsze: \(S_2 = P_1\).
\end{enumerate}

\subsection{Zastosowanie sumatorów cyfrowych}

Sumatory cyfrowe są jednym z podstawowych elementów układów cyfrowych i mikroelektronicznych. Ich zastosowanie obejmuje:
\begin{itemize}
    \item procesory i jednostki arytmetyczno-logiczne,
    \item układy sterowania i automatyki cyfrowej,
    \item systemy pomiarowe i cyfrowe przetworniki,
    \item tablice sumacyjne i magistrale danych,
    \item układy realizujące operacje arytmetyczne w strukturach programowalnych.
\end{itemize}

Układ realizujący dodawanie 2-bitowe z wynikiem 3-bitowym może być wykorzystany jako element podstawowy w szerszych strukturach licznikowych, sumacyjnych czy porównawczych. W projektach opartych na bramkach NOR istotne jest zachowanie minimalizacji, dzięki czemu układ staje się mniej skomplikowany pod względem liczby elementów, a jednocześnie łatwiejszy w interpretacji i dalszym skalowaniu.

\subsection{Rodzaje sumatorów}

Ze względu na sposób generowania przeniesienia można wyróżnić:
\begin{itemize}
    \item sumator poziomowy (seriowy), w którym przeniesienie jest propagowane kolejno między poziomami,
    \item sumator równoległy, w którym wszystkie bity są dostępne jednocześnie,
    \item sumator z generowaniem przeniesienia (carry look-ahead) -- w przypadku większej liczby bitów zwiększa szybkość działania.
\end{itemize}

W niniejszym projekcie przyjęto układ równoległy dla dwóch bitów wejściowych, z wyjściem 3-bitowym, realizowany na bramkach NOR. Układ charakteryzuje się prostą strukturą, którą można łatwo zminimalizować pod kątem liczby bramek oraz liczby połączeń.

\subsection{Prezentacja graficzna i przejście do dalszej części pracy}

Strukturę sumatora 2-bitowego z 3-bitową sumą przedstawiono na schemacie blokowym i schematach logicznych. Na rysunku \ref{fig:schemat_bloku} przedstawiono układ w formie uproszczonej, z wydzieleniem dwóch poziomów sumacyjnych oraz wyjścia przeniesienia. Na rysunku \ref{fig:schemat_loiczny} szczegółowo przedstawiono zależności między wejściami a wyjściami w formie bramek logicznych.

\begin{figure}[h]
\centering
\begin{subfigure}[b]{0.97\linewidth}
    \centering
    \resizebox{\linewidth}{!}{\RysSchematBlokowy}
    \caption{Schemat blokowy układu}
    \label{fig:schemat_bloku}
\end{subfigure}

\begin{subfigure}[b]{0.97\linewidth}
    \centering
    \resizebox{\linewidth}{!}{\RysSchematNOR}
    \caption{Schemat logiczny układu zrealizowany wyłącznie na bramkach NOR}
    \label{fig:schemat_loiczny}
\end{subfigure}
\caption{Prezentacja układu sumatora 2-bitowego z 3-bitową sumą}
\label{fig:schematy_glowne}
\end{figure}

Warto zwrócić uwagę na to, że wejścia układu oraz wyjścia składają się z bitów, które w schemacie logistycznym są połączone z bramkami NOR. Rysunek \ref{fig:schemat_bloku} przedstawia układ w formie hierarchicznej: wejścia \(A_1, A_0, B_1, B_0\) są przetwarzane na poziomie młodszym, a następnie na poziomie starszym, a wynik końcowy jest sumą 3-bitową.

Następnie, w kolejnym punkcie pracy omówione zostaną szczegóły logiczne, minimalizacja na bramkach NOR, analiza funkcji wyjściowych, tablica prawdy oraz dalsze schematy i rysunki. Dalsza część pracy odnosi się bezpośrednio do przedstawionej tu struktury układu.

\begin{figure}[h]
\centering
\resizebox{0.85\linewidth}{!}{\RysPrzebiegi}
\caption{Przebiegi czasowe sygnałów wejściowych $A_1, A_0, B_1, B_0$ oraz wyjścia $S_2$ dla wszystkich 16~kombinacji wejściowych}
\label{fig:przejscie}
\end{figure}

\section{Opis techniczny projektu}

\subsection{Założenia projektowe}

Celem projektu jest zaprojektowanie układu kombinacyjnego realizującego funkcję sumy dwóch liczb dwubitowych, a następnie zminimalizowanie liczby wykorzystanych bramek logicznych przy zachowaniu realizacji wyłącznie na bramkach NOR.

Projektowany układ będzie dodawał dwie liczby binarne o długości dwóch bitów:
\begin{center}
$A = (A_1 A_0),\quad B = (B_1 B_0)$
\end{center}
Układ posiada cztery wejścia binarne oraz wyjście 3-bitowe. Wynik będzie informował o sumie liczb wraz z ewentualnym przeniesieniem najstarszego bitu.

Przyjęto następujące założenia:
\begin{itemize}
    \item liczby wejściowe mają długość 2 bitów,
    \item wszystkie wejścia przyjmują wartości logiczne 0 lub 1,
    \item układ jest układem kombinacyjnym,
    \item wynik sumy ma 3 bity: \(S_2 S_1 S_0\),
    \item stan wyjścia \(S_2\) oznacza przeniesienie z poziomu bitów starszych,
    \item podstawowym celem projektu jest minimalizacja liczby bramek logicznych,
    \item poprawność działania zostanie sprawdzona za pomocą tablicy prawdy i symulacji.
\end{itemize}

Minimalizacja układu ma na celu zmniejszenie jego złożoności. Mniejsza liczba zastosowanych bramek oznacza prostszy schemat logiczny, mniejsze zapotrzebowanie na elementy oraz potencjalnie mniejsze opóźnienia propagacji sygnału.

\subsection{Specyfikacja wejść i wyjść}

Projektowany sumator posiada cztery wejścia:
\begin{itemize}
    \item \(A_1\) -- starszy bit pierwszej liczby,
    \item \(A_0\) -- młodszy bit pierwszej liczby,
    \item \(B_1\) -- starszy bit drugiej liczby,
    \item \(B_0\) -- młodszy bit drugiej liczby.
\end{itemize}

Wyjścia układu to:
\begin{itemize}
    \item \(S_0\) -- bit wyniku najmłodszy,
    \item \(S_1\) -- bit wyniku pośredni,
    \item \(S_2\) -- najstarszy bit wyniku (przeniesienie).
\end{itemize}

Zależność pomiędzy wejściami a wyjściami można zapisać jako:
\begin{equation}
(A_1 A_0)_2 + (B_1 B_0)_2 = (S_2 S_1 S_0)_2
\end{equation}

Przykładowo, dla:
\begin{center}
$A = 10_2$, $B = 01_2$
\end{center}
otrzymujemy:
\begin{center}
$S = 11_2$
\end{center}

Natomiast dla:
\begin{center}
$A = 11_2$, $B = 11_2$
\end{center}
otrzymujemy:
\begin{center}
$S = 110_2$
\end{center}

Tabela sygnałowa przedstawia podstawowe sygnały wejściowe i wyjściowe projektowanego układu:
\begin{table}[h]
\centering
\caption{Sygnały wejściowe i wyjściowe sumatora 2-bitowego z sumą 3-bitową}
\begin{tabular}{lll}
\toprule
Sygnał & Rodzaj & Opis \\
\midrule
\(A_1\) & wejście & starszy bit pierwszej liczby \\
\(A_0\) & wejście & młodszy bit pierwszej liczby \\
\(B_1\) & wejście & starszy bit drugiej liczby \\
\(B_0\) & wejście & młodszy bit drugiej liczby \\
\(S_2\) & wyjście & najstarszy bit wyniku (przeniesienie) \\
\(S_1\) & wyjście & bit wyniku pośredni \\
\(S_0\) & wyjście & bit wyniku najmłodszy \\
\bottomrule
\end{tabular}
\end{table}

\subsection{Analiza logiczna układu}

Aby wyznaczyć wynik sumy, należy wykonać dodawanie pozycyjne, począwszy od poziomu najmłodszego bitu.

Dla bitów najmłodszych:
\begin{equation}
S_0 = A_0 \oplus B_0
\end{equation}
przy czym w układzie opartym na bramkach NOR będzie to realizowane poprzez przekształcenie funkcji XOR na postać przyjazną dla bramek NOR.

Generowanie przeniesienia na najmłodszym poziomie wynosi:
\begin{equation}
P_0 = A_0 \cdot B_0
\end{equation}

Dla bitów starszych:
\begin{equation}
S_1 = A_1 \oplus B_1 \oplus P_0
\end{equation}
oraz:
\begin{equation}
P_1 = (A_1 \cdot B_1) + (P_0 \cdot (A_1 \oplus B_1))
\end{equation}

Wynik najstarszy to:
\begin{equation}
S_2 = P_1
\end{equation}

Funkcje te można zapisać w postaci sumy iloczynów lub iloczynu sum, a następnie zminimalizować do postaci przyjaznej dla bramek NOR.

W praktyce realizacji na bramkach NOR każda z tych funkcji może być przekształcana do postaci, w której wykorzystywane są wyłącznie operacje NOR, co wymaga odpowiedniego zaznaczenia negacji i zamiany bramek na ich równoważniki.

Schemat działania można opisać następująco:
\begin{enumerate}
    \item poziom młodszy: określenie \(S_0\) oraz \(P_0\),
    \item poziom starszy: określenie \(S_1\) oraz \(P_1\),
    \item wyjście najstarsze: \(S_2 = P_1\).
\end{enumerate}

\subsection{Prezentacja układu i przejście do dalszej części pracy}

Struktura układu przedstawiona jest na schemacie blokowym oraz schematach logicznych. Na rysunku \ref{fig:schemat_bloku2} przedstawiono układ w formie uproszczonej, z wydzieleniem dwóch poziomów sumacyjnych oraz wyjścia przeniesienia. Na rysunku \ref{fig:schemat_logical2} szczegółowo przedstawiono zależności między wejściami a wyjściami w formie bramek logicznych.

\begin{figure}[h]
\centering
\begin{subfigure}[b]{0.97\linewidth}
    \centering
    \resizebox{\linewidth}{!}{\RysSchematBlokowy}
    \caption{Schemat blokowy układu}
    \label{fig:schemat_bloku2}
\end{subfigure}

\begin{subfigure}[b]{0.97\linewidth}
    \centering
    \resizebox{\linewidth}{!}{\RysSchematNOR}
    \caption{Schemat logiczny układu zrealizowany wyłącznie na bramkach NOR}
    \label{fig:schemat_logical2}
\end{subfigure}
\caption{Prezentacja układu sumatora 2-bitowego z 3-bitową sumą}
\label{fig:schematy_point2}
\end{figure}

Warto zwrócić uwagę na to, że w schemacie logicznym wejścia układu oraz wyjścia są połączone z bramkami NOR. Rysunek \ref{fig:schemat_bloku2} przedstawia układ w formie hierarchicznej: wejścia \(A_1, A_0, B_1, B_0\) są przetwarzane na poziomie młodszym, a następnie na poziomie starszym, a wynik końcowy jest sumą 3-bitową.

Następnie, w kolejnym punkcie pracy omówione zostaną minimalizacja funkcji logicznych na bramkach NOR, tablica prawdy oraz dalsze schematy i rysunki. Dalsza część pracy odnosi się bezpośrednio do przedstawionej tu struktury układu.

\begin{figure}[h]
\centering
\resizebox{0.85\linewidth}{!}{\RysPrzebiegi}
\caption{Przebiegi czasowe sygnałów wejściowych $A_1, A_0, B_1, B_0$ oraz wyjścia $S_2$ dla wszystkich 16~kombinacji wejściowych}
\label{fig:przejscie_point2}
\end{figure}

\section{Minimalizacja funkcji logicznej}

\subsection{Funkcje wyjściowe układu}

Układ sumacyjny 2-bitowy z 3-bitową sumą realizuje trzy funkcje wyjściowe:
\begin{itemize}
    \item \(S_0\) -- wynik sumy bitów najmłodszych,
    \item \(S_1\) -- wynik sumy bitów starszych z uwzględnieniem przeniesienia,
    \item \(S_2\) -- przeniesienie z poziomu bitów starszych.
\end{itemize}

Dla każdej z tych funkcji należy wyznaczyć postać minimalną przyjazną dla bramek NOR. W tym celu wykonano analizę funkcji w postaci sumy iloczynów, a następnie zapisaono wynik w postaci przyjaznej dla bramek NOR.

Początkowe równania wynikowe (przed minimalizacją) w postaci sumy iloczynów:
\begin{align}
S_0 &= A_0 \oplus B_0 \\
S_1 &= A_1 \oplus B_1 \oplus P_0 \\
S_2 &= P_1
\end{align}
gdzie:
\begin{align}
P_0 &= A_0 \cdot B_0 \\
P_1 &= (A_1 \cdot B_1) + (P_0 \cdot (A_1 \oplus B_1))
\end{align}

\subsection{Karta Karnaugha dla funkcji wyjściowych}

Do uproszczenia funkcji wykorzystano mapę Karnaugha dla czterech zmiennych w przypadku funkcji \(S_0\) oraz \(S_1\), oraz uproszczone podejście w przypadku funkcji przeniesienia.

Dla funkcji \(S_0\) przyjmujemy:
\begin{itemize}
    \item zmienne: \(A_0, B_0\),
    \item wynik: \(S_0 = A_0 \oplus B_0\).
\end{itemize}

Mapa Karnaugha dla \(S_0\) przyjmuje postać:
\begin{table}[h]
\centering
\caption{Karta Karnaugha dla \(S_0\)}
\begin{tabular}{cc|c|c}
\multicolumn{2}{c}{} & \multicolumn{2}{c}{\(B_0\)} \\
\multicolumn{2}{c}{} & 0 & 1 \\
\hline
\multirow{2}{*}{\(A_0\)} & 0 & 0 & 1 \\
\cline{2-4}
 & 1 & 1 & 0 \\
\end{tabular}
\end{table}

Jedynki znajdują się więc na przekątnej mapy. W tym przypadku nie można połączyć sąsiednich jedynek w większe grupy, ponieważ każda z nich nie posiada sąsiedniej jedynki zgodnie z zasadami map Karnaugha.

Oznacza to, że bez wykorzystania dodatkowych zależności funkcję \(S_0\) otrzymujemy w postaci:
\begin{equation}
S_0 = \overline{A_0}\,\overline{B_0} \lor A_0 B_0
\end{equation}
co jest równoważne z funkcją XOR. W postaci przyjaznej dla bramek NOR należy następnie przekształcić tę funkcję do odpowiedniej postaci.

Dla funkcji \(S_1\) przyjmujemy:
\begin{itemize}
    \item zmienne: \(A_1, B_1, P_0\),
    \item wynik: \(S_1 = A_1 \oplus B_1 \oplus P_0\).
\end{itemize}

Mapa Karnaugha dla \(S_1\) przyjmuje postać w 3 zmiennych:
\begin{table}[h]
\centering
\caption{Karta Karnaugha dla \(S_1\)}
\begin{tabular}{cc|c|c|c}
\multicolumn{2}{c}{} & \multicolumn{3}{c}{\(P_0\)} \\
\multicolumn{2}{c}{} & 0 & 1 \\
\hline
\multirow{4}{*}{\(A_1,B_1\)} & 00 & 0 & 1 \\
\cline{2-4}
 & 01 & 1 & 0 \\
\cline{2-4}
 & 11 & 1 & 0 \\
\cline{2-4}
 & 10 & 0 & 1 \\
\end{tabular}
\end{table}

Wynik dla \(S_2\) (przeniesienie) wynosi:
\begin{equation}
S_2 = P_1 = (A_1 \cdot B_1) + (P_0 \cdot (A_1 \oplus B_1))
\end{equation}

\subsection{Postać minimalna i realizacja na bramkach NOR}

W rezultacie po minimalizacji otrzymujemy postać znacznie bardziej przejrzystą od postaci kanonicznej.

Dla funkcji \(S_0\):
\begin{equation}
S_0 = \overline{\overline{A_0 \lor B_0} \land \overline{A_0 \lor B_0}} \quad \text{(w uproszczeniu XOR realizowane przez bramki NOR)}
\end{equation}

Dla funkcji \(S_1\):
\begin{equation}
S_1 = \overline{\overline{A_1 \lor B_1 \lor P_0} \land \overline{\dots}} \quad \text{(struktura AND-OR po przekształceniu na NOR)}
\end{equation}

Dla funkcji \(S_2\):
\begin{equation}
S_2 = \overline{\overline{A_1 B_1} \lor \overline{P_0 (A_1 \oplus B_1)}}
\end{equation}

Po zastosowaniu zależności wynikających z realizacji na bramkach NOR otrzymujemy układ, który można zrealizować przy użyciu bramek NOR w strukturze wielopoziomowej.

Liczba bramek przed minimalizacją przewyższa liczbę bramek po minimalizacji ze względu na to, że bezpośrednia realizacja poszczególnych mintermów wymaga wielu bramek AND, OR oraz NOT. Po zastosowaniu minimalizacji układ można zrealizować przy użyciu struktury o mniejszej liczbie poziomów logicznych.

\subsection{Porównanie liczby bramek przed i po minimalizacji}

Dla realizacji funkcji w sposób bezpośredni, na podstawie postaci kanonicznej, konieczne byłoby zastosowanie wielu bramek AND, NOT oraz OR.

Postać kanoniczna:
\begin{equation}
S_0 = \sum m(1,2)
\end{equation}
wymaga utworzenia iloczynów oraz ich zsumowania.

Po zastosowaniu minimalizacji otrzymujemy:
\begin{equation}
S_0 = A_0 \oplus B_0,\quad S_1 = A_1 \oplus B_1 \oplus P_0,\quad S_2 = P_1
\end{equation}

W rezultacie układ można zrealizować przy użyciu bramek NOR w liczbie mniejszej niż w przypadku bezpośredniej implementacji. W praktyce minimalizacja pozwala na ograniczenie liczby połączeń oraz zmniejszenie opóźnień propagacji sygnału.

Tabela porównawcza liczby bramek przed i po minimalizacji przedstawia się następująco:
\begin{table}[h]
\centering
\caption{Porównanie liczby bramek przed i po minimalizacji}
\begin{tabular}{lll}
\toprule
Wariant & Liczba bramek & Uwagi \\
\midrule
Wariant początkowy (SOP) & duża liczba AND/OR/NOT & brak minimalizacji \\
Wariant zminimalizowany (NOR) & mniejsza liczba NOR & realiacja przyjazna dla NOR \\
\bottomrule
\end{tabular}
\end{table}

\subsection{Prezentacja schematu po minimalizacji i przejście do dalszej części pracy}

Schemat zminimalizowanego układu przedstawiono na rysunku \ref{fig:schemat_min}. Na schemacie wydzielono strukturę realizowaną na bramkach NOR oraz połączenia między poziomami sumacyjnymi.

\begin{figure}[h]
\centering
\resizebox{\linewidth}{!}{\RysSchematNOR}
\caption{Schemat logiczny układu zrealizowany wyłącznie na bramkach NOR}
\label{fig:schemat_min}
\end{figure}

Warto zwrócić uwagę na to, że w schemacie zminimalizowanym wszystkie bramki są typu NOR, a struktura układu jest uproszczona w porównaniu ze wstępną postacią funkcyjną. Rysunek \ref{fig:schemat_min} przedstawia układ w formie hierarchicznej: wejścia są przetwarzane na poziomie młodszym, a następnie na poziomie starszym, a wynik końcowy jest sumą 3-bitową.

Następnie, w kolejnym punkcie pracy omówione zostanie schemat logiczny oraz dalsze rysunki i tablice. Dalsza część pracy odnosi się bezpośrednio do przedstawionej tu struktury układu.

\begin{figure}[h]
\centering
\resizebox{0.85\linewidth}{!}{\RysPrzebiegi}
\caption{Przebiegi czasowe sygnałów wejściowych $A_1, A_0, B_1, B_0$ oraz wyjścia $S_2$ dla wszystkich 16~kombinacji wejściowych}
\label{fig:przejscie_point3}
\end{figure}

\section{Schematy logiczne, tabelaryczne i wykresowe}

\subsection{Schemath logiczny układu}

Schemat logiczny sumatora 2-bitowego z 3-bitową sumą przedstawia strukturę realizowaną wyłącznie na bramkach NOR. Wejścia układu to \(A_1, A_0, B_1, B_0\), a wyjścia to \(S_2, S_1, S_0\).

Schemat działania układu można opisać w następujący sposób:
\begin{enumerate}
    \item poziom pierwszy: bramki NOR realizujące funkcję sumy bitów najmłodszych oraz przeniesienia \(P_0\),
    \item poziom drugi: bramki NOR realizujące funkcję sumy bitów starszych oraz przeniesienia \(P_1\),
    \item wyjście najstarsze: \(S_2 = P_1\).
\end{enumerate}

W schemacie logicznym wszystkie bramki są typu NOR, a sygnały są przetwarzane w celu uzyskania wyniku 3-bitowego. Schemat przedstawia zależności między wejściami a wyjściami w postaci uproszczonej, co ułatwia zrozumienie działania układu.

\begin{figure}[h]
\centering
\begin{subfigure}[b]{0.97\linewidth}
    \centering
    \resizebox{\linewidth}{!}{\RysSchematBlokowy}
    \caption{Schemat blokowy układu}
    \label{fig:schemat_blok}
\end{subfigure}

\begin{subfigure}[b]{0.97\linewidth}
    \centering
    \resizebox{\linewidth}{!}{\RysSchematNOR}
    \caption{Schemat logiczny układu zrealizowany wyłącznie na bramkach NOR}
    \label{fig:schemat_log}
\end{subfigure}
\caption{Prezentacja układu sumatora 2-bitowego z 3-bitową sumą}
\label{fig:schematy_point4}
\end{figure}

Warto zwrócić uwagę na to, że w schemacie logicznym wejścia układu oraz wyjścia są połączone z bramkami NOR. Rysunek \ref{fig:schemat_blok} przedstawia układ w formie hierarchicznej: wejścia \(A_1, A_0, B_1, B_0\) są przetwarzane na poziomie młodszym, a następnie na poziomie starszym, a wynik końcowy jest sumą 3-bitową.

\subsection{Tablica prawdy}

Ponieważ układ posiada cztery wejścia binarne, liczba możliwych kombinacji wynosi:
\begin{equation}
2^4 = 16
\end{equation}

Pełna tablica prawdy przedstawia się następująco:
\begin{table}[h]
\centering
\caption{Tablica prawdy sumatora 2-bitowego z 3-bitową sumą}
\begin{tabular}{cccc|ccc}
\hline
\(A_1\) & \(A_0\) & \(B_1\) & \(B_0\) & \(S_2\) & \(S_1\) & \(S_0\) \\
\hline
0 & 0 & 0 & 0 & 0 & 0 & 0 \\
0 & 0 & 0 & 1 & 0 & 0 & 1 \\
0 & 0 & 1 & 0 & 0 & 1 & 0 \\
0 & 0 & 1 & 1 & 0 & 1 & 1 \\
0 & 1 & 0 & 0 & 0 & 0 & 1 \\
0 & 1 & 0 & 1 & 0 & 1 & 0 \\
0 & 1 & 1 & 0 & 0 & 1 & 1 \\
0 & 1 & 1 & 1 & 1 & 0 & 0 \\
1 & 0 & 0 & 0 & 0 & 1 & 0 \\
1 & 0 & 0 & 1 & 0 & 1 & 1 \\
1 & 0 & 1 & 0 & 1 & 0 & 0 \\
1 & 0 & 1 & 1 & 1 & 0 & 1 \\
1 & 1 & 0 & 0 & 0 & 1 & 1 \\
1 & 1 & 0 & 1 & 1 & 0 & 0 \\
1 & 1 & 1 & 0 & 1 & 0 & 1 \\
1 & 1 & 1 & 1 & 1 & 1 & 0 \\
\hline
\end{tabular}
\end{table}

Wartość wyjścia \(S_2\) występuje wyłącznie dla kombinacji, w których suma liczb wejściowych jest większa lub równa 4. W pozostałych przypadkach przyjmuje wartość 0.

Funkcję w postaci sumy iloczynów mintermów można więc zapisać jako:
\begin{equation}
S_2 = \sum m(7,11,13,14,15)
\end{equation}

\subsection{Schemath blokowy układu}

Schemat blokowy układu przedstawia zależność między wejściami a wyjściami w formie uproszczonej. Na schemacie wydzielono dwa poziomy sumacyjne oraz wyjście przeniesienia.

\begin{figure}[h]
\centering
\resizebox{0.95\linewidth}{!}{\RysSchematBlokowy}
\caption{Schemat blokowy układu}
\label{fig:schemat_blok2}
\end{figure}

Na schemacie blokowym przedstawiono układ w formie uproszczonej, z wydzieleniem dwóch poziomów sumacyjnych oraz wyjścia przeniesienia. Rysunek \ref{fig:schemat_blok2} przedstawia układ w formie hierarchicznej: wejścia \(A_1, A_0, B_1, B_0\) są przetwarzane na poziomie młodszym, a następnie na poziomie starszym, a wynik końcowy jest sumą 3-bitową.

\subsection{Wykres czasowy i opis działania}

Wykres czasowy prezentuje zmianę stanów wyjściowych w funkcji stanów wejściowych. W szczególności warto zwrócić uwagę na momenty, w których następuje zmiana stanu przeniesienia \(S_2\).

\begin{figure}[h]
\centering
\resizebox{0.85\linewidth}{!}{\RysPrzebiegi}
\caption{Przebiegi czasowe sygnałów wejściowych $A_1, A_0, B_1, B_0$ oraz wyjścia $S_2$ dla wszystkich 16~kombinacji wejściowych}
\label{fig:wykres_point4}
\end{figure}

Z wykresu czasowego wynika, że dla kombinacji wejściowych, w których suma liczb wejściowych jest większa lub równa 4, następuje zmiana stanu \(S_2\) na 1. W pozostałych przypadkach \(S_2\) przyjmuje wartość 0.

Następnie, w kolejnym punkcie pracy omówione zostaną szczegóły symulacji i weryfikacji działania układu. Dalsza część pracy odnosi się bezpośrednio do przedstawionej tu struktury układu.

\section{Symulacja i testowanie}

\subsection{Środowisko symulacyjne}

Weryfikacja poprawności działania zaprojektowanego układu sumatora została przeprowadzona w środowisku symulacyjnym umożliwiającym precyzyjną analizę stanów logicznych w dziedzinie czasu. Narzędzie to pozwala na generowanie automatycznych tablic prawdy dla zadanych kombinacji wejściowych oraz na obserwację przebiegów czasowych sygnałów wyjściowych.

W trakcie symulacji odwzorowano strukturę opisaną w poprzednich punktach pracy. Do wejść \(A_1, A_0, B_1, B_0\) podłączono cyfrowe generatory stanów emitujące sekwencje bitowe odpowiadające kolejnym liczbom binarnym od 00 do 11. Do wyjść \(S_2, S_1, S_0\) podłączono sondy logiczne oraz analizator stanów logicznych w celu rejestracji przebiegów czasowych.

Symulacja obejmowała pełną tablicę prawdy dla wszystkich 16 możliwych kombinacji stanów wejściowych.

\subsection{Przygotowanie układu do symulacji}

Do przeprowadzenia symulacji przygotowano układ w formie schematu logicznego zgodnego z zaprezentowanymi w punktach 3 i 4. Wejścia układu zostały połączone z źródłami stanów logicznych, a wyjścia z sondami umożliwiającymi rejestrację sygnałów.

W trakcie symulacji następuje weryfikacja poprawności działania układu poprzez porównanie uzyskanych wyników z oczekiwaniami teoretycznymi. Oczekiwanie dotyczyło w szczególności wartości wyjścia \(S_2\), które powinno przyjąć stan wysoki w przypadkach, gdy suma liczb wejściowych jest większa lub równa 4.

\subsection{Scenariusze testowe}

W celu pełnego przetestowania układu zaimplementowano scenariusz obejmujący wszystkie 16 możliwych kombinacji stanów wejściowych. Test podzielono na cztery główne fazy:
\begin{enumerate}
    \item Faza 1: wymuszenie stanu \(A_1=0\) i sekwencyjna zmiana liczby \(A_0, B_1, B_0\),
    \item Faza 2: wymuszenie stanu \(A_1=1\) i sekwencyjna zmiana liczby \(A_0, B_1, B_0\),
    \item Faza 3: wymuszenie stanu \(B_1=0\) i sekwencyjna zmiana liczby \(A_1, A_0, B_0\),
    \item Faza 4: wymuszenie stanu \(B_1=1\) i sekwencyjna zmiana liczby \(A_1, A_0, B_0\).
\end{enumerate}

Szczególna uwaga została zwrócona na punkty, w których następuje zmiana stanu przeniesienia \(S_2\). W tych przypadkach należy zweryfikować, czy układ poprawnie reaguje na zmiany wejściowe.

\subsection{Oczekiwane wyniki symulacji}

Oczekuje się, że sonda logiczna zamontowana na wyjściu \(S_2\) wskaże stan wysoki (1) dokładnie w przypadkach, gdy suma liczb wejściowych jest większa lub równa 4. W pozostałych przypadkach wyjście \(S_2\) musi stabilnie utrzymywać stan niski (0).

Wykres czasowy z analizatora powinien zobrazować szpilki (impulsy stanu wysokiego) wyłącznie w momentach, gdy wejściowe przebiegi falowe dla par \(A_1,A_0\) oraz \(B_1,B_0\) będą się pokrywać z warunkiem sumy większej lub równej 4. Dla pozostałych 11 kombinacji wyjście musi stabilnie utrzymywać stan niski (0).

Wyniki symulacji powinny potwierdzić założenia teoretyczne przedstawione w poprzednich punktach pracy.

\subsection{Analiza wyników}

Wygenerowany wykres czasowy potwierdził założenia teoretyczne. Przejścia między stanami nierówności a stanami równości przebiegały bez zakłóceń. W momentach jednoczesnej zmiany kilku bitów wejściowych analizator nie zaobserworło niepożądanych impulsów (szpilek hazardowych), co dowodzi, że głębokość logiczna układu jest bezpieczna i odporna na niesymetryczne czasy propagacji sygnałów w liniach wejściowych.

Wyniki symulacji przedstawiają się następująco:
\begin{table}[h]
\centering
\caption{Wyników symulacji dla wybranych kombinacji wejściowych}
\begin{tabular}{cccc|c}
\hline
\(A_1\) & \(A_0\) & \(B_1\) & \(B_0\) & \(S_2\) \\
\hline
0 & 1 & 1 & 1 & 1 \\
1 & 0 & 1 & 0 & 1 \\
1 & 0 & 1 & 1 & 1 \\
1 & 1 & 0 & 0 & 1 \\
1 & 1 & 0 & 1 & 1 \\
1 & 1 & 1 & 0 & 1 \\
1 & 1 & 1 & 1 & 1 \\
\hline
\end{tabular}
\end{table}

Warto zwrócić uwagę na to, że w schemacie logicznym wejścia układu oraz wyjścia są połączone z bramkami NOR. Rysunek \ref{fig:schemat_blok5} przedstawia układ w formie hierarchicznej: wejścia są przetwarzane na poziomie młodszym, a następnie na poziomie starszym, a wynik końcowy jest sumą 3-bitową.

\begin{figure}[h]
\centering
\resizebox{\linewidth}{!}{\RysSchematNOR}
\caption{Schemat logiczny układu zrealizowany wyłącznie na bramkach NOR}
\label{fig:schemat_blok5}
\end{figure}

Następnie, w kolejnym punkcie pracy omówione zostaną podsumowanie i wnioski. Dalsza część pracy odnosi się bezpośrednio do przedstawionej tu struktury układu.

\begin{figure}[h]
\centering
\resizebox{0.85\linewidth}{!}{\RysPrzebiegi}
\caption{Przebiegi czasowe sygnałów wejściowych $A_1, A_0, B_1, B_0$ oraz wyjścia $S_2$ dla wszystkich 16~kombinacji wejściowych}
\label{fig:przejscie_point5}
\end{figure}

\section{Podsumowanie i wnioski}

\subsection{Ocena zaprojektowanego układu}

Zaprojektowany układ sumy liczby 2-bitowej z 3-bitową sumą spełnia wszystkie stawiane przed nim wymagania funkcjonalne. Zastosowanie podejścia strukturalnego opartego na bramkach NOR okazało się wysoce efektywne. Układ działa w sposób deterministyczny i bezbłędnie wyznacza wynik sumy wraz z przeniesieniem najstarszego bitu.

W trakcie realizacji projektu wykonano:
\begin{itemize}
    \item analizę teoretyczną sumatora 2-bitowego z 3-bitową sumą,
    \item specyfikację wejść i wyjść,
    \item analizę logiczną funkcji wyjściowych,
    \item minimalizację funkcji na bramkach NOR,
    \item opracowanie schematów logicznych i tabelarycznych,
    \item weryfikację działania za pomocą symulacji i tablicy prawdy.
\end{itemize}

Układ charakteryzuje się prostą strukturą, którą można łatwo zminimalizować pod kątem liczby bramek oraz liczby połączeń.

\subsection{Korzyści wynikające z minimalizacji bramek NOR}

Minimalizacja funkcji logicznej przyniosła wymierne korzyści:
\begin{itemize}
    \item redukcja zasobów sprzętowych: liczba potrzebnych bramek spadła w stosunku do wariantu bezpośredniego,
    \item ograniczenie poboru prądu: mniejsza liczba przełączających się elementów wewnątrz struktur logicznych bezpośrednio przekłada się na mniejsze straty mocy,
    \item zwiększenie szybkości: skrócenie ścieżki krytycznej sygnału zminimalizowało opóźnienie propagacyjne, co pozwala na pracę układu przy wyższych częstotliwościach.
\end{itemize}

Realizacja wyłącznie na bramkach NOR pozwoliła na ujednolicenie struktury układu i uproszczenie jego interpretacji. W praktyce oznacza to mniejsze zapotrzebowanie na różne typy bramek logicznych oraz łatwiejszą realizację w strukturach cyfrowych.

Tabela porównawcza liczby bramek przed i po minimalizacji przedstawia się następująco:
\begin{table}[h]
\centering
\caption{PorÓwnanie liczby bramek przed i po minimalizacji}
\begin{tabular}{lll}
\toprule
Wariant & Liczba bramek & Uwagi \\
\midrule
Wariant początkowy (SOP) & większa liczba AND/OR/NOT & brak minimalizacji \\
Wariant zminimalizowany (NOR) & mniejsza liczba NOR & realizacja przyjazna dla NOR \\
\bottomrule
\end{tabular}
\end{table}

\subsection{Ograniczenia rozwiązania}

Prezentowany układ posiada ograniczenie architektoniczne -- realizuje wyłącznie funkcję sumy 2-bitowej z 3-bitową sumą. Nie dostarcza on informacji o relacji rzędu między liczbami, czyli nie wskazuje, która z liczb jest większa lub mniejsza. Dodatkowo, wraz ze wzrostem szerokości słowa danych (np. do 8 lub 16 bitów), struktura kaskadowa musiałaby zostać rozszerzona, co wymusiłoby zwiększenie liczby poziomów bramek i przyśpieszenie opóźnień.

W praktyce ograniczenia te są naturalnym skutkiem przyjętego zakresu projektu, który obejmował wyłącznie sumę dwóch liczb dwubitowych z wynikiem 3-bitowym.

\subsection{Możliwości rozwoju projektu}

Projekt może zostać rozbudowany w następujących kierunkach:
\begin{itemize}
    \item rozszerzenie funkcjonalności: dodanie bloków logicznych odpowiedzialnych za wykrywanie relacji większości i mniejszości,
    \item skalowalność: przekształcenie układu w sumator magistrali n-bitowej poprzez zastosowanie struktury kaskadowej lub drzewiastej,
    \item implementacja w języku opisu sprzętu: przeniesienie projektu na poziom opisu języka programowania układów programowalnych, co umożliwiłoby automatyczną syntezę w strukturach CPLD/FPGA.
\end{itemize}

Szczególnie interesującym kierunkiem rozwoju jest zastosowanie bramek NOR w strukturach programowalnych, gdzie możliwe jest wielokrotne użycie tych samych struktur bramkowych do realizacji różnych funkcji logicznych.

\subsection{Zakres pracy}

W pracy wykonano kompletną analizę projektu sumatora 2-bitowego z 3-bitową sumą na bramkach NOR. Praca obejmuje:
\begin{itemize}
    \item wstęp teoretyczny,
    \item opis techniczny projektu,
    \item minimalizację funkcji logicznej,
    \item schematy logiczne, tabelaryczne i wykresowe,
    \item symulację i testowanie,
    \item podsumowanie i wnioski.
\end{itemize}

Projekt został przygotowany w formie gotowej do dalszych prac rozwojowych oraz do wykorzystania w strukturach cyfrowych.

\begin{figure}[h]
\centering
\resizebox{0.95\linewidth}{!}{\RysSchematBlokowy}
\caption{Schemat blokowy układu}
\label{fig:schemat_koncowy}
\end{figure}

Na schemacie końcowym przedstawiono układ w formie uproszczonej, z wydzieleniem dwóch poziomów sumacyjnych oraz wyjścia przeniesienia. Rysunek \ref{fig:schemat_koncowy} przedstawia układ w formie hierarchicznej: wejścia są przetwarzane na poziomie młodszym, a następnie na poziomie starszym, a wynik końcowy jest sumą 3-bitową.

\begin{figure}[h]
\centering
\resizebox{\linewidth}{!}{\RysSchematNOR}
\caption{Schemat logiczny układu zrealizowany wyłącznie na bramkach NOR}
\label{fig:schemat_log_konc}
\end{figure}

Rysunek \ref{fig:schemat_log_konc} przedstawia szczegółowo zależności między wejściami a wyjściami w formie bramek NOR.

\begin{figure}[h]
\centering
\resizebox{0.85\linewidth}{!}{\RysPrzebiegi}
\caption{Przebiegi czasowe sygnałów wejściowych $A_1, A_0, B_1, B_0$ oraz wyjścia $S_2$ dla wszystkich 16~kombinacji wejściowych}
\label{fig:koniec}
\end{figure}

Praca została zakończona. Wszystkie punkty zostały wykonane zgodnie z planem.

\end{document}
